\documentclass[]{article}

\usepackage{amsmath}
\usepackage{multirow}
\usepackage{natbib}
\usepackage{graphicx} 


\begin{document}

Anticipating shifts in systemic risk is a foundational challenge in financial economics, with the CBOE Volatility Index (VIX) serving as the principal barometer of expected market-wide turmoil. While many forecasting models rely on aggregate market data, valuable precursory information may lie latent within the market's internal structure. The dynamics of how different economic sectors behave relative to one another can offer insight into the underlying stability of the financial system, raising a critical question: can the divergence in the risk profiles of individual sectors provide a leading signal for future changes in aggregate market volatility?

The central problem this paper addresses is the ambiguity inherent in measures of cross-sectional volatility dispersion. An increase in the dispersion of realized volatility across sectors—where some parts of the market become significantly more turbulent than others—can be interpreted in two distinct ways. It could represent benign, idiosyncratic noise confined to specific industries, or it could be a symptom of rising market fragility, where underlying economic stressors affect sectors heterogeneously before coalescing into a systemic event. Distinguishing between these two interpretations is crucial for developing reliable risk-monitoring tools that can filter meaningful signals from market noise.

To resolve this ambiguity, we propose and test a conditional framework. We first construct a daily metric, termed Sectoral Volatility Dispersion (SVD), calculated as the cross-sectional standard deviation of 21-day realized volatilities across ten major US sector ETFs. Our central hypothesis is that the predictive power of SVD for future VIX innovations is not absolute but is contingent on the market's broader correlational structure. Specifically, we posit that elevated SVD only becomes a potent indicator of impending systemic risk when it is accompanied by a breakdown in the market's internal cohesion, measured by the average cross-sector correlation. This approach frames market fragility not just as a function of rising idiosyncratic risk, but as the simultaneous occurrence of such risk and the decoupling of its constituent parts.

Our empirical investigation, which employs a Hidden Markov Model to endogenously classify VIX regimes, provides strong support for this conditional relationship. We find that SVD, when considered in isolation, is not a statistically significant predictor of future VIX innovations or the probability of transitioning into a high-volatility state. However, its predictive power emerges robustly through its interaction with market structure. Elevated SVD is significantly associated with higher 21-day ahead VIX innovations, but only when average cross-sector correlation is low. This study contributes by clarifying that cross-sectional volatility dispersion should not be interpreted as a standalone timing signal. Instead, it is a critical component of a more sophisticated indicator of market fragility, where the combination of high idiosyncratic volatility and sector decoupling signals a heightened vulnerability to future market-wide stress.

\end{document}
                